% whatsnew.tex -- reader's guide to the updated report. \newfigcount and \newtabcount come from
% sections/_counts.tex, which tools/assemble.py generates by counting the figure/table environments of Part II.
\newcommand{\newfigcount}{77}
\newcommand{\newtabcount}{24}


\section*{What is new in this version}
\addcontentsline{toc}{section}{What is new in this version}

The Checkpoint~2 report was compiled on 18~September 2026. This version brings it level with the project
repository\footnote{\url{https://github.com/Abhi-Manyu12/htn-biology}, branch \texttt{main}, commit
\texttt{61d18d8} of 12~November 2026, which was still the head of the branch on 17~November 2026.} by
comparing the report against every package, result file and document the repository contains.
The table below is a guide to what was added. Part~\ref{part:new} contains \newfigcount{} figures
and \newtabcount{} tables that Checkpoint~2 did not have.

\begin{table}[H]
\centering
\small
\caption*{Guide to this version. Sections, algorithms, tables and figures of Checkpoint~2 keep their original numbers.}
\begin{tabularx}{\linewidth}{@{}p{0.30\linewidth}X@{}}
\toprule
Section & What it contains \\
\midrule
\multicolumn{2}{@{}l}{\textit{Part~\ref{part:cp2}: the Checkpoint 2 report}} \\
Sections~\ref{sec:orig-prelim}--\ref{sec:orig-summary}
  & The Checkpoint~2 text, unchanged, with 15 dated update notes where the repository now shows that a
    statement is out of date or was incorrect. \\
\addlinespace
\multicolumn{2}{@{}l}{\textit{Part~\ref{part:new}: work since Checkpoint 2}} \\
\ref{sec:record}\quad Project record
  & What was built and when; the structure of the repository; the roadmap status table of Checkpoint~2
    brought up to date; related work and positioning. \\
\ref{sec:core}\quad Core HTN model
  & The three plans that exist and which one Checkpoint~2 showed; spiral geometry; the anchor-removal sweep to
    80\,\% and the wind sweep; the baseline comparison re-examined; the simulated webs of the interactive
    viewer; a table of corrections. \\
\ref{sec:astar_fap}\quad A* and fixed action pattern
  & Two formulations Checkpoint~2 does not mention, measured and compared descriptively with the HTN plans. \\
\ref{sec:mulder}\quad Moving supports
  & A benchmark after Mulder et al.: six execution policies building a web while its supports move; pilot,
    controller audit and validation studies; galleries of the simulated webs. \\
\ref{sec:recovery}\quad Re-entrant recovery
  & Recovery from failures during execution with IPyHOP; solution trees before and after repair; the audit
    showing that the planner as configured does not backtrack. \\
\ref{sec:sim}\quad Physical simulation
  & A two-dimensional simulation of the planned web in which threads sag under gravity and respond to wind,
    moving supports and anchor loss, with its validation against the analytic catenary. \\
\ref{sec:record-b}\quad Claims and limitations
  & The repository's ledger of 28 claims with the status of each; how results are kept reproducible; what
    remains open. \\
\ref{sec:synthesis}\quad Revised discussion
  & The conclusions of Checkpoint~2 set against the present evidence; issues found while preparing this
    update; conclusions. \\
\bottomrule
\end{tabularx}
\end{table}

\paragraph{How the figures were made.}
Every plotted value comes from a file committed in the repository or from a re-run of the repository's own
code at a stated commit, and each caption names its source. The scripts that produced the figures, and a
manifest listing the source data of each one, are kept with the source of this report under \texttt{scripts/} and
\texttt{figures/}. Drawings of simulated webs are either the repository's own renderings, converted
unchanged, or redrawn from the plan and trace data it stores.

\paragraph{How to read the results.}
The repository keeps a ledger of what may and may not be claimed (Section~\ref{sec:record-claims}), and this
report follows it. Several of the new results are negative. They are reported in the same detail as the
positive ones, because they change what the earlier conclusions can bear
(Table~\ref{tab:synthesis-revisited}).
